\subsection{为什么 inner 需要多步：从精确 KL 几何到 Natural Gradient}
\label{sec:ch2-natural-gradient}

前面的 mirror 解是在策略分布空间中定义的，而实际训练变量是神经网络参数 $\theta$。因此需要把“固定一个 outer target 后的 inner 优化”从分布空间搬到参数空间。下面先不假设局部二次模型，而是从 softmax 参数化和 KL 的定义逐步展开。

固定一个 prefix $s$，令 $z_\theta(s)\in\mathbb R^{|\mathcal A|}$ 为学生 logits，$\pi_\theta(\cdot\mid s)=\operatorname{softmax}(z_\theta(s))$。在第 $k$ 个 outer step 中冻结 $\pi_k$ 和 $r_k$，把 inner 的镜像目标写成最小化形式
\begin{equation}
\label{eq:ch2-opd-inner-objective}
\mathcal J_k(\theta)
=
-\E_{s,a}[r_k(s,a)]
+\frac{1}{\eta}\E_s\!\left[
\KL\!\left(\pi_\theta(\cdot\mid s)\Vert\pi_k(\cdot\mid s)\right)
\right],
\end{equation}
其中 $\E_{s,a}$ 的 action 测度可以是 $\pi_\theta$、$\pi_k$ 或 replay 加权测度；下文的局部曲率只要求该测度固定且具有有限二阶矩。

\paragraph{第一步：logit 空间的 KL 一阶项消失。}
令 $\delta z=z_\theta-z_{\theta_k}$，其中 $\pi_{\theta_k}=\pi_k$。对固定的行为分布 $q=\pi_k$ 定义
\[
K(z)=\KL(\operatorname{softmax}(z)\Vert q).
\]
直接对 softmax 的 log-partition 展开：
\[
K(z_k+\delta z)
=
\sum_a p_a\log p_a-\sum_a p_a\log q_a,
\qquad p=\operatorname{softmax}(z_k+\delta z).
\]
因为 $p=\pi_k=q$ 时 KL 取全局最小值，$\nabla_zK(z_k)=0$。进一步使用
\[
\frac{\partial p_a}{\partial z_b}
=p_a(\mathbf 1\{a=b\}-p_b)
\]
可得 Hessian
\begin{equation}
\label{eq:ch2-logit-fisher}
\nabla_z^2K(z_k)
=
F_k(s)
=
\operatorname{Diag}(\pi_k)-\pi_k\pi_k^{\mathsf T}.
\end{equation}
$F_k$ 是 categorical policy 的 Fisher 矩阵，半正定且零空间由所有常数 logit 平移张成。因此
\begin{equation}
\label{eq:ch2-kl-logit-expansion}
\KL(\pi_\theta\Vert\pi_k)
=
\frac12\delta z^{\mathsf T}F_k\delta z
+R_{\mathrm{KL}}(\delta z),
\qquad
|R_{\mathrm{KL}}(\delta z)|
\leq C_{\mathrm{KL}}\|\delta z\|^3
\end{equation}
在 logits 位于有界集合时成立。这个三阶余项来自 softmax 的三阶导数有界，而不是额外的二元 KL 假设。

\paragraph{第二步：reward 项的一阶展开。}
令
\[
\ell_k(\theta)
=
-\E_{s,a}[r_k(s,a)]
\]
表示最小化目标中的 reward 部分。假设 $\ell_k$ 在 $\theta_k$ 邻域内二阶可微，则
\begin{equation}
\label{eq:ch2-reward-taylor}
\ell_k(\theta_k+\Delta\theta)
=
\ell_k(\theta_k)
-b_k^{\mathsf T}\Delta\theta
+R_r(\Delta\theta),
\qquad
b_k=\E_{s,a}\!\left[
r_k(s,a)\nabla_\theta\log\pi_{\theta_k}(a\mid s)
\right],
\end{equation}
并且 $|R_r(\Delta\theta)|\leq C_r\|\Delta\theta\|^2$。这里 $b_k$ 是固定 outer target 在当前学生处的 policy-gradient 方向；如果使用 PMD-Part 或 BPO 的回归 loss，则 $b_k$ 被相应的 regression score 替代，但 KL 曲率项相同。

\paragraph{第三步：从 logits Jacobian 得到参数空间曲率。}
令
\[
J_k(s)=\left.\frac{\partial z_\theta(s)}{\partial\theta}\right|_{\theta=\theta_k}.
\]
由
\[
\delta z
=
J_k(s)\Delta\theta+R_z(\Delta\theta),
\qquad
\|R_z(\Delta\theta)\|\leq C_z\|\Delta\theta\|^2,
\]
把式 \eqref{eq:ch2-kl-logit-expansion} 代入式 \eqref{eq:ch2-opd-inner-objective}，得到
\begin{equation}
\label{eq:ch2-kl-parameter-expansion}
\E_s\KL(\pi_\theta\Vert\pi_k)
=
\frac12\Delta\theta^{\mathsf T}A_k\Delta\theta
+R_{\mathrm{par}}(\Delta\theta),
\qquad
A_k=\E_s[J_k(s)^{\mathsf T}F_k(s)J_k(s)],
\end{equation}
其中 $|R_{\mathrm{par}}(\Delta\theta)|\leq C\|\Delta\theta\|^3$。因此式 \eqref{eq:ch2-local-quadratic-derived} 不是直接假设，而是式 \eqref{eq:ch2-kl-parameter-expansion} 与式 \eqref{eq:ch2-reward-taylor} 合并后，在 $\Delta\theta$ 足够小时的二阶主项：
\begin{equation}
\label{eq:ch2-local-quadratic-derived}
\mathcal J_k(\theta_k+\Delta\theta)
=
\mathcal J_k(\theta_k)
+\frac{1}{2\eta}\Delta\theta^{\mathsf T}A_k\Delta\theta
-b_k^{\mathsf T}\Delta\theta
+R_k(\Delta\theta).
\end{equation}

\begin{proposition}[局部曲率的正半定性与不可辨识方向]
\label{prop:ch2-local-curvature}
$A_k$ 对称半正定。若某个参数方向 $v$ 满足 $J_k(s)v$ 对所有 replay prefix 都是常数 logit 平移，则 $A_kv=0$；该方向不改变策略分布，因而 KL 不能为它提供曲率。若在商空间上 $A_k\succeq mI$，则镜像 inner 目标在该商空间局部强凸。
\end{proposition}

\begin{proof}
对任意 $v$，
\[
v^{\mathsf T}A_kv
=
\E_s[(J_k(s)v)^{\mathsf T}F_k(s)(J_k(s)v)]\geq0
\]
由 $F_k\succeq0$ 成立。$F_k$ 的零空间正是常数向量方向，所以 $J_k(s)v$ 为常数时对应二次型为零。去除这些不改变策略的参数方向后，若最小特征值至少为 $m>0$，则二次项给出局部强凸性。
\end{proof}

\subsection*{有限 inner 步的收敛条件}

定义局部主模型
\[
Q_k(\Delta\theta)
=
\frac{1}{2\eta}\Delta\theta^{\mathsf T}A_k\Delta\theta
-b_k^{\mathsf T}\Delta\theta.
\]
若 $b_k$ 属于 $A_k$ 的值域，则其最小范数解为
\[
\Delta\theta_k^\star=\eta A_k^\dagger b_k.
\]
使用普通梯度下降、步长 $\gamma>0$：
\begin{equation}
\label{eq:ch2-inner-gd-derived}
\Delta\theta_{m+1}
=
\Delta\theta_m-\gamma\nabla Q_k(\Delta\theta_m)
=
(I-\gamma A_k/\eta)\Delta\theta_m+\gamma b_k.
\end{equation}

\begin{theorem}[由 KL 曲率推出的 inner 线性收敛]
\label{thm:ch2-local-gd-derived}
设 $A_k$ 在商空间上的特征值属于
$[m_k,M_k]$，且 $0<m_k\leq M_k$。若
\[
0<\gamma<\frac{2\eta}{M_k},
\]
则式 \eqref{eq:ch2-inner-gd-derived} 收敛到 $\Delta\theta_k^\star$，并满足
\[
\|\Delta\theta_m-\Delta\theta_k^\star\|_{A_k}
\leq
q_k^m\|\Delta\theta_0-\Delta\theta_k^\star\|_{A_k},
\qquad
q_k=\max_{\lambda\in[m_k,M_k]}
\left|1-\frac{\gamma\lambda}{\eta}\right|<1.
\]
特别地，取 $\gamma=2\eta/(m_k+M_k)$ 时，
\[
q_k=\frac{M_k-m_k}{M_k+m_k}
=\frac{\kappa_k-1}{\kappa_k+1},
\qquad
\kappa_k=M_k/m_k.
\]
\end{theorem}

\begin{proof}
令 $e_m=\Delta\theta_m-\Delta\theta_k^\star$。由式 \eqref{eq:ch2-inner-gd-derived}，
\[
e_{m+1}=(I-\gamma A_k/\eta)e_m.
\]
在 $A_k$ 的特征方向 $\lambda$ 上，误差乘以 $1-\gamma\lambda/\eta$。给定步长条件，所有这些因子的绝对值都小于 $1$，因此得到 $q_k^m$ 的界。令两个端点方向的绝对收缩率相等即可得到最优固定步长和最后的条件数表达式。
\end{proof}

真实 inner 目标还包含余项 $R_k$。若其梯度在半径 $r$ 的邻域满足
\[
\|\nabla R_k(\Delta\theta)\|
\leq L_k\|\Delta\theta\|^2,
\]
并且 $r$ 足够小使得 $L_k r\leq m_k/(2\eta)$，则真实 Hessian 的特征值仍落在
\[
[m_k/(2\eta),\,M_k/\eta+L_kr]
\]
内。于是标准强凸光滑函数的梯度下降结论仍成立，只是收缩率由上述扩大后的光谱区间决定。这给出了“inner 需要做到局部收敛”的明确条件：步长必须由 KL/Fisher 曲率控制，且迭代必须停留在二阶展开有效的信赖域内。

\subsection*{Natural Gradient 的来源与作用}

普通梯度在参数坐标中使用 $A_k$ 的谱，病态方向的收敛速度由 $\kappa_k$ 决定。Natural Gradient 以 Fisher 度量约束一次更新：
\begin{equation}
\label{eq:ch2-natural-gradient-subproblem}
\min_{\Delta\theta}\ 
-b_k^{\mathsf T}\Delta\theta
\quad\text{s.t.}\quad
\frac12\Delta\theta^{\mathsf T}A_k\Delta\theta\leq\frac{\eta}{2}.
\end{equation}
拉格朗日一阶条件给出
\[
\Delta\theta_{\mathrm{NG}}
=
c\,A_k^\dagger b_k,
\]
其中常数 $c$ 由 trust-region 半径确定。若直接取 $\Delta\theta=\eta A_k^\dagger b_k$，它正是局部二次模型的最小点；实际算法用 Fisher 近似 $\widehat A_k$、阻尼 $\lambda I$ 和有限 inner 步，因此得到的是近似自然梯度，而不是理论上无条件的一步收敛。

\begin{proposition}[Natural Gradient 对参数重参数化的不变性]
\label{prop:ch2-ng-invariance}
设 $\vartheta=\phi(\theta)$ 是局部可逆的光滑重参数化，且 Fisher 矩阵按张量变换规则变换。则更新方向
\[
\Fisher_\theta^\dagger\nabla_\theta\mathcal J_k
\]
与在 $\vartheta$ 坐标中计算后映射回 $\theta$ 的方向一致；普通梯度 $\nabla_\theta\mathcal J_k$ 一般不具备该性质。
\end{proposition}

\begin{proof}
令 $B=\partial\vartheta/\partial\theta$。梯度变换为 $\nabla_\theta\mathcal J=B^{\mathsf T}\nabla_\vartheta\mathcal J$，Fisher 变换为 $\Fisher_\theta=B^{\mathsf T}\Fisher_\vartheta B$。在可逆子空间上代入并使用伪逆的相应变换，得到
\[
\Fisher_\theta^\dagger\nabla_\theta\mathcal J
=
B^{-1}\Fisher_\vartheta^\dagger\nabla_\vartheta\mathcal J.
\]
普通梯度只含 $B^{\mathsf T}$，因此会随坐标尺度改变。
\end{proof}

这解释了 MDPO 和 EffOPD 中的两个实现判断。只在 $\theta_k$ 处做一次梯度步时，KL 的一阶梯度为零，trust-region 曲率尚未进入更新；固定 $\pi_k$ 并执行多次 inner 更新，才能逼近式 \eqref{eq:ch2-opd-inner-objective}。Natural Gradient 进一步把步长按照 $A_k$ 的几何尺度归一化，但它仍然需要信赖域、阻尼和有限 inner 预算来处理模型非线性与 Fisher 估计误差。


\paragraph{几何中间分布与算术 proximal teacher 的区别。} 式 \eqref{eq:ch2-opd-geometric-bridge} 给出的 $\pi_k^{1-\eta}T^\eta/Z_k$ 是 KL 镜像子问题的精确解，因此在 log-probability 空间插值。TOP-D 中的算术 teacher
\[
\widetilde T_\eta=(1-\eta)\pi_k+\eta T
\]
是另一个用于控制 log-ratio 方差的 teacher operator。它同样产生 student 与 teacher 之间的中间分布，但不等价于 KL mirror 解；前者改变 target teacher，后者改变 proximal 优化问题的解。
